\crefalias{section}{appendix}
\section{Derivation of Hypercharge (Geometric Flux Conservation)}
\label{app:OriginOfHypercharge}

In the foundational derivations of the $E_8$-Persistence substrate \cite{meyer_informational_2026}, the structure of the non-Abelian gauge groups $SU(3)_C$ and $SU(2)_L$ emerged strictly from the geometric invariants $\sigma=5$ and $\chi=2$. In standard Quantum Field Theory, the corresponding $U(1)_Y$ hypercharge assignments ($+1/3, +4/3, -2/3, -1, -2$) are axiomatic inputs, justified post-hoc by the requirement of anomaly cancellation.

In the $E_8$-Persistence theory, anomaly cancellation is \textbf{Information Conservation (Unitarity)}. Information (phase flux) cannot be created or destroyed at a lattice node. By applying this strict geometric flux conservation to the $\nu=16$ chiral channels of the substrate, we can derive the exact fractional hypercharge assignments entirely from first principles.

\subsection{Step 1: The Channel Partition (\texorpdfstring{$\nu=16$}{nu=16})}
The lattice node possesses exactly $\nu=16$ chiral channels. These channels are geometrically partitioned by the topological boundary ($\chi=2$) and the internal bulk ($\sigma-\chi=3$):

\begin{itemize}
    \item \textbf{Color Partition:} Channels either couple to the internal interaction bulk (Color, multiplicity 3) or they bypass it (Colorless, multiplicity 1). 
    \item \textbf{Boundary Partition:} Channels either couple to the active temporal boundary (Left-Handed Doublets, multiplicity 2) or they remain static internal singlets (Right-Handed, multiplicity 1).
\end{itemize}
This strict geometric intersection partitions the $\nu=16$ capacity into exactly the Standard Model fermion architecture: 
\begin{equation}
(3 \times 2) + (3 \times 1) + (3 \times 1) + (1 \times 2) + (1 \times 1) + (1 \times 1) = \mathbf{16}
\end{equation}
This perfectly generates the quark doublets ($Q_L$), the right-handed quarks ($u_R, d_R$), the lepton doublet ($L_L$), the right-handed electron ($e_R$), and one sterile channel ($\nu_R$), utilizing 100\% of the node's hardware capacity.

\subsection{Step 2: The Fractional Denominator (Color Dilution)}
Hypercharge ($U(1)_Y$) represents a continuous phase rotation across the node. To preserve the $SU(3)$ interaction symmetry, any flux entering the color bulk must be distributed isotropically across its dimensions. 

Because the color bulk possesses exactly $\sigma-\chi=3$ geometric degrees of freedom, the total flux for a color-coupled state must be divided by 3. Therefore, quarks inherently possess fractional hypercharges with a denominator of $\mathbf{3}$, while color-singlet states (leptons) couple directly to the manifold and carry integer charges.

\subsection{Step 3: Geometric Flux Conservation}
To preserve Unitarity, the net chiral phase flux (Left-Handed input minus Right-Handed output) must balance to zero across every orthogonal geometric sector.

\paragraph{Boundary Flux Balance:}
The total flux traversing the $\chi=2$ topological boundary must sum to zero. The three color-coupled doublets ($Y_{Q}$) must perfectly balance the single colorless doublet ($Y_L$).
\begin{equation}
3 Y_Q + 1 Y_L = 0
\end{equation}
Because the denominator is strictly fixed to 3 by color dilution, the minimal positive integer unit of flux available to this boundary state is exactly $+1/3$. Assigning this fundamental unit ($Y_Q = +1/3$) strictly forces the lepton doublet flux: $Y_L = -1$. Geometrically, the Lepton doublet exactly balances the combined flux of the three Quark doublets.

\paragraph{Internal Bulk Balance:}
The chiral flux entering the color bulk from the boundary doublets ($2 \times Y_Q$) must equal the flux returning from the right-handed color singlets ($Y_u + Y_d$).
\begin{equation}
Y_u + Y_d = 2 Y_Q = \frac{2}{3}
\end{equation}
How does the lattice geometrically split this $2/3$ flux between the two right-handed quark states? Because they are singlets under the boundary, their values must be defined by the remaining orthogonal invariants of the projection: the spatial manifold ($D=4$) and the boundary reflection ($-\chi = -2$). The constraint $Y_u + Y_d = 2/3$ combined with the requirement that numerators be geometric invariants of the projection ($\{1, \chi, D, \sigma\}$) admits a unique solution: $Y_u = D/3 = 4/3$ and $Y_d = -\chi/3 = -2/3$ perfectly satisfies this conservation law:
\begin{equation}
Y_u + Y_d = \frac{D - \chi}{3} = \frac{4 - 2}{3} = \frac{2}{3}
\end{equation}

\paragraph{Total Manifold Balance:}
Finally, the net chiral flux across all active channels on the macroscopic manifold must be zero ($\sum_{LH} Y - \sum_{RH} Y = 0$). We established that the Left-Handed boundary sum is zero. Therefore, the Right-Handed output sum must also be zero:
\begin{equation}
3 Y_u + 3 Y_d + 1 Y_e = 0
\end{equation}
Substituting the derived quark fluxes ($3(4/3) + 3(-2/3) = 4 - 2 = 2$), we solve for the right-handed electron:
\begin{equation}
2 + Y_e = 0 \implies Y_e = -2
\end{equation}
Geometrically, $-2$ is exactly $-\chi$. The right-handed electron carries the exact inverse flux of the topological boundary.

\subsection{Conclusion: The Derived Charge Structure}
This geometric framework naturally recovers the hypercharge assignments without requiring supplementary anomaly cancellation conditions. By strictly enforcing information conservation (Unitarity) across the $\nu=16$ chiral channels, the lattice geometry mathematically forces the exact fractional charges of the Standard Model. 

Table \ref{tab:hypercharge_derivation} summarizes this derivation, demonstrating that every fermion's hypercharge is identically the geometric ratio required to conserve flux between the manifold ($D=4$), the color bulk ($\sigma-\chi=3$), and the topological boundary ($\chi=2$).

\begin{table}[ht!]
\centering
\caption{Derived Fermion Hypercharges via Geometric Flux Conservation}
\label{tab:hypercharge_derivation}
\renewcommand{\arraystretch}{1.2}
\begin{tabular}{@{}lccccc@{}}
\toprule
\textbf{Particle} & \textbf{Lattice Partition} & \textbf{Multiplicity} & \textbf{Derived $Y$} & \textbf{Geometric Ratio} \\
\midrule
Quark ($Q_L$)    & Boundary \& Color & 3 & $+1/3$ & $+1/(\sigma-\chi)$ \\
Lepton ($L_L$)   & Boundary & 1 & $-1$ & $-3 \times Y_Q$ \\
\addlinespace
Up Quark ($u_R$)    & Manifold \& Color & 3 & $+4/3$ & $+D/(\sigma-\chi)$ \\
Down Quark ($d_R$)  & Reflection \& Color & 3 & $-2/3$ & $-\chi/(\sigma-\chi)$ \\
Electron ($e_R$)    & Reflection & 1 & $-2$ & $-\chi$ \\
\bottomrule
\end{tabular}
\end{table}

\subsection{Derivation of Higgs Hypercharge}
As established in System 5, the scalar sector acts as the regulatory bridge anchored to the lattice Unitary Ground State ($\Delta^0 = 1$). Geometrically, this mandates a unitary hypercharge for the scalar doublet: $Y_H = +1$. 

This purely geometric assignment perfectly recovers the observed physical charge states of the Higgs field. Evaluating the Gell-Mann--Nishijima relation ($Q = I_3 + Y/2$) with this geometric hypercharge ($Y_H = +1$) and the $SU(2)_L$ isospin components ($I_3 = \pm 1/2$) yields exactly the neutral ($H^0$, $Q=0$) and charged ($H^+$, $Q=+1$) components of the Standard Model Higgs doublet.